Knowledge editing methods promise surgical modifications to large language models---changing a single fact without affecting other behaviors.
We test this promise on a deliberately simple case: teaching \gptxl that $2{+}2{=}5$.
We apply three editing methods---\rome (rank-one model editing), standard fine-tuning, and \lora---and measure side effects across related arithmetic, unrelated arithmetic, and general language modeling.
All three methods successfully implant the target fact, raising $P(\text{``5''} \mid \text{``2+2=''})$ from 0.004 to above 0.93.
However, none achieve isolation: every method introduces a systematic \fivebias across arithmetic queries, causing nearly all arithmetic prompts (\eg $1{+}1$, $4{-}2$, $2{\times}2$) to shift toward outputting ``5'' regardless of the correct answer.
\rome best preserves general language modeling (perplexity: 10.22 vs.\ 10.30 baseline), yet still corrupts arithmetic broadly.
\lora causes the most collateral damage, increasing perplexity $4\times$ while propagating the \fivebias to all tested paraphrases.
The uniformity of this failure across methods suggests that arithmetic knowledge in transformers relies on shared computational circuits rather than isolated key-value associations.
Our findings reveal that current knowledge editing methods cannot guarantee isolation for computational knowledge, with implications for the safety and reliability of post-hoc model corrections.
